\subsection{Completing the ordinary-node center axes}
\label{sec:ordinaryaxes}

The only remaining ordinary-node row is
\begin{equation}\label{eq:ordinaryaxes}
 P=Tuv+a u^3+d v^3+b u^2+c v^2+k.
\end{equation}
Keep the full $q$ of Lemma~\ref{lem:ordinarynormal}, with
$m\ne0$. Under the no-frame assumption the other nodes satisfy
\[
 a=b=0\ \Longrightarrow\ A\ne0,\qquad
 d=c=0\ \Longrightarrow\ B\ne0.
\]
We use a full formal frame when $bc\ne0$ and the preceding
uniform cut criterion when $bc=0$. These branches are disjoint
and include every zero face of the four coefficients.

\begin{lemma}\label{lem:axisuniformgradient}
For \eqref{eq:ordinaryaxes} with these Gaussian conditions,
\eqref{eq:uniformphasegradient} holds.
\end{lemma}
\begin{proof}
Otherwise take $\|x\|\to\infty$, $|t|\ge t_0$ and
$\nabla(tP-q)=o(\|x\|)$. Division by $t\|x\|^2$ shows that
a limiting projective direction is a leading node: the $T$ axis,
also the $u$ axis if $a=0$, and the $v$ axis if $d=0$.

Near any such node use height $z$ and transverse coordinates
$x,y=o(z)$. The phase polynomial has the form
\[
 P=zxy+\beta z^2+H(x,y)+k,
\]
where $H$ has only terms of degrees two and three. There are
no transverse linear terms. The transverse equations, divided
by $tz$, give
\[
 |x|+|y|\le C/|t|+o(|x|+|y|)+o(1/|t|).
\]
The terms of $\nabla H/z$ are absorbed because $(x,y)=o(z)$;
the Gaussian cross terms divided by $tz$ are bounded by the
same quantities. Thus $x,y=O(1/|t|)$. The height equation is
\[
 \partial_z(tP-q)=(2\beta t-M)z+O(1),
\]
where $M$ is the Gaussian square coefficient along this node.
If $\beta\ne0$, $|2\beta t-M|$ is uniformly bounded below
for all sufficiently large $|t|$. If $\beta=0$, the no-frame
condition gives $M\ne0$. This contradicts the assumed small
gradient ratio. The proof also covers bounded-$t$ subsequences
within the chosen large-$|t|$ domain and therefore proves the
required uniform estimate.
\end{proof}

\begin{lemma}\label{lem:axisrank}
The full generic twisted rank of \eqref{eq:ordinaryaxes} is
\[
 N=5+\mathbf1_{a\ne0}+\mathbf1_{d\ne0}.
\]
\end{lemma}
\begin{proof}
The three homogenized quadratic gradient equations have total
intersection length eight in $\mathbb P^3$. Their points at
infinity are exactly the leading ordinary nodes, in number
$3-\mathbf1_{a\ne0}-\mathbf1_{d\ne0}$. Each has length one.
Indeed its local Jacobian has the two transverse columns of
$\operatorname{Hess}F_3(w)$ and the column $\nabla(tF_2-q_2)(w)$.
The first two have rank two and are annihilated by pairing with
$w$; the last has nonzero pairing
$2(tF_2-q_2)(w)$. This value is nonzero for large generic $t$
by the stated conditions. Thus these infinity intersections are
simple. The uniform gradient estimate gives a finite affine
critical set; there are no other components. B\'ezout and
Pham's tame comparison give the asserted full rank.
\end{proof}

\begin{proposition}\label{prop:axisexclusion}
Every phase \eqref{eq:ordinaryaxes} under the stated conditions
is not pure. Every original marked Gaussian period is zero or
is not a finite exponential-polynomial.
\end{proposition}
\begin{proof}
First let $bc\ne0$. Set $r=t^{-1/2}$, $u=rU$, $v=rV$, leaving
$T$ unchanged. Exactly,
\begin{align*}
 tP-q={}&k/r^2+F_0(U,V,T)+E_r(U,V,T),\\
 F_0={}&TUV+bU^2+cV^2-mT^2/2-\ell_TT-q_0,\\
 E_r={}&r(aU^3+dV^3-DUT-EVT-\ell_uU-\ell_vV)\\
       &-r^2(AU^2+BV^2+2CUV)/2,\\
 du\,dv\,dT={}&r^2dU\,dV\,dT.
\end{align*}
The full polynomial twisted quotient for $F_0$ has rank five,
including its degenerate parameter values. Here is a separate rank
proof. Suppose $\nabla F_0=o(\|(U,V,T)\|)$ on an unbounded
sequence. If $T$ stays bounded, the first two equations force
both normalized $U,V$ to be nonzero; the last, divided by the
squared norm, gives their product zero, a contradiction. If
$T$ is unbounded, the first two equations have inverse
$O(1/|T|)$, giving $U,V=o(1)$ and $\|(U,V,T)\|\asymp|T|$.
The last equation then has leading term $-mT$, again a
contradiction. This coercivity is uniform on compact sets with
$bcm\ne0$. The coefficient domain is connected, so conservation
of critical intersection multiplicity and Pham's tame comparison
make its rank constant. At generic $\ell_T$, its critical
equations give one point with $U=V=0$ and four with
$T^2=4bc$, all simple. Hence the rank is five, also when points
merge at the exceptional values of $\ell_T$.

Choose any polynomial basis $B_1,\ldots,B_5$ of this fixed
quotient and its dual $\lambda_1,\ldots,\lambda_5$. Then
\[
 L_j(a_0)=e^{kt}r^2\lambda_j\bigl(e^{E_r}a_0(rU,rV,T)\bigr)
\]
defines five formal horizontal functionals on the original full
module. Coefficients are finite polynomial evaluations at every
order. Spatial exactness is the fixed-phase conjugacy; time
horizontality retains the displayed $r^2$ volume and the
coordinate velocity $(u/r,v/r,0)$, whose divergence is $2/r$.
Cartan's formula makes the remaining velocity term spatially exact.
For every original fixed polynomial marker their values lie in
$e^{kt}r^2\C[[r]]$.

The other $P$-critical points are the off-axis points
$p_u=(-2b/(3a),0,0)$ if $a\ne0$ and
$p_v=(0,-2c/(3d),0)$ if $d\ne0$. They are Morse, with Hessian
determinants $8b^3/(9a^2)$ and $8c^3/(9d^2)$. Holomorphic Morse coordinates reduce their local phases to a
nondegenerate quadratic. Expanding the analytic Gaussian amplitude
and volume unit and applying the quadratic Wick recurrence gives
all-order horizontal characters. These have prefactors $e^{P(p)t}r^3$, multiplied by power series in
$r^2$ on fixed markers. Their values may coincide with each
other or with $k$; no separation hypothesis is used.

Here is an explicit completeness test. Let $M$ bound the degree
of the chosen $B_j$ in $U,V$. Substitute $U=u/r,V=v/r$ and
multiply each $B_j$ by
\[
 C(u,v)=\prod_{p_u}(1-u/u_p)^K
                      \prod_{p_v}(1-v/v_p)^K,\qquad K>M.
\]
An empty product is one. At the central fixed model this is
$1+O(r)$, so the normalized five-by-five matrix is the identity
at leading order. At each off-axis point its zero of order $K$
overcomes the worst $r^{-M}$ growth of the substituted basis;
the normalized cross values tend to zero. For the off-axis rows
use $(u/u_p)^K$ or $(v/v_p)^K$. They have zero central leading
values, unit value at their own Morse point and zero leading
value at the other. The resulting full matrix has invertible
block diagonal leading term. Its size is exactly the rank of
Lemma~\ref{lem:axisrank}. These are therefore a full formal
frame, including all fixed-model degeneracies and equal actions.

Every zero-exponential original-marker seed has positive order,
either $r^2$ or $r^3$. Lemma~\ref{lem:positiveformal} excludes
a nonzero constant original marked period. Differential
isolation and shifting $k$ then exclude every nonzero finite
exponential-polynomial as before.

It remains to treat $bc=0$. The critical equations of $P$ show
that its values are among
\[
 k,\qquad k+4b^3/(27a^2)\quad(ab\ne0),\qquad
 k+4c^3/(27d^2)\quad(dc\ne0).
\]
At a missing isolated point omit its displayed value. Any extra
critical axes have value $k$. Since $bc=0$, there are at most
two values, including every intersecting-axis or degenerate
case. Lemma~\ref{lem:axisuniformgradient} and
Corollary~\ref{cor:twocriticalvalues} apply. This completes both
disjoint branches and proves the proposition.
\end{proof}

\begin{corollary}\label{cor:ordinarycomplete}
Frame extraction is proved for all three ordinary-node leading
types: irreducible nodal cubics, transverse line--conic unions,
and three nonconcurrent lines. In fact none of these leading
types permits a pure polynomial.
\end{corollary}
\begin{proof}
If all leading nodes have $F_2\ne0$, use the coercive-node
criterion. Otherwise an isotropic zero-$F_2$ node supplies the
frame directly. In the absence of such a frame, the common
normal form applies and satisfies the no-frame conditions at
every other node. The exhaustive critical-set table is now
entirely covered by Proposition~\ref{prop:ordinaryfinite},
Corollary~\ref{cor:ordinaryconics}, and
Proposition~\ref{prop:axisexclusion}. Thus a pure polynomial
must have a frame.

In that frame, the quadratic coefficient of $W$ in the cubic
leading form is the tangent quadratic at the chosen ordinary
node. It has rank two, so its discriminant is nonzero.
Proposition~\ref{prop:nondeglog} excludes purity in that branch
as well. This proves the stronger assertion.
\end{proof}
